\section{INTRODUCTION}
\label{sec:introduction}

\subsection{Background and Motivation}

Communication and sensing are often needed within the same system. A mobile
terminal needs a data link and knowledge of its position, while a network
operator needs to carry traffic and monitor the infrastructure that supports
it. Integrated sensing and communication (ISAC) brings these functions together
through shared signals, hardware, or infrastructure
\cite{Zhang2026ISACEvolution}. This interest extends to future mobile networks,
with ISAC included among the usage scenarios in the International Mobile
Telecommunications framework for 2030 and beyond
\cite{ITUR_M2160_2023}.

Optical integrated sensing and communication (O-ISAC) applies this idea to
optical links and photonic systems. A fiber can carry communication and sensing
signals along the same route, while a visible light source can support data
transmission and positioning \cite{OISAC_SCR00013,OISAC_SCR00196}.
Figure~\ref{fig:oisac_platform_overview} illustrates representative ways of
combining the two functions. The role of optics differs across these systems.
In fiber and optical wireless links, light propagates through the sensing
medium or reaches the receiver or target. In photonics-enabled wireless
systems, optical components can generate the signal while propagation to the
target takes place in the millimeter-wave or terahertz band
\cite{OISAC_SCR00036,OISAC_SCR00083}.

\begin{figure*}[!t]
\centering
\includegraphics[width=\textwidth]{figures/fig_introduction_oisac_platforms.png}
\caption{Representative O-ISAC configurations in fiber, free-space optical,
visible-light, and photonics-enabled terahertz systems. Each panel shows a
shared resource and the communication and sensing functions it supports.
The design choices below the panels can affect both functions, with their
roles depending on the implementation. These examples illustrate physical
configurations rather than a performance ranking or an exhaustive taxonomy.}
\label{fig:oisac_platform_overview}
\end{figure*}

Sharing a resource can also link the performance of the two functions.
In one fiber experiment, raising the sensing-probe power while holding
communication launch power fixed improved the sensing signal-to-noise ratio
but increased the communication bit error rate \cite{OISAC_SCR00057}.
Joint operation can also allow sensing to assist communication. In a
visible-light prototype, position estimates guided beamforming toward the
receiver and improved the communication data rate at the tested locations
\cite{OISAC_SCR00196}. These examples show why the shared resource and the way
it is used matter. Joint design must account for both the benefit sought and
its effect on the other function.

The physical setting determines how these relationships can be used.
A vibration measurement along a fiber and a position estimate in a room
answer different sensing questions. Optical wireless links also depend on
propagation and alignment, while fiber systems have their own transmission
and sensing constraints \cite{OISAC_SCR00916,OISAC_SCR01085,Mohsan2026OISACReview}.
A researcher selecting an O-ISAC approach therefore needs more than a list
of achievable rates or sensing errors. The useful questions are which design
choice produced the result, how both functions responded, and whether the
conditions match the intended application. This connects the reported
performance to decisions about how to build and evaluate a system.

\subsection{Related Surveys}
\label{sec:related_surveys}

Existing reviews provide both platform-specific design guidance and broader
views of O-ISAC. Wen et al.\ \cite{Wen2024OISACArchitectures} explain optical
architectures and waveform allocation, using a common simulation setting to
show communication error rate and ranging error. Liang et al.
\cite{Liang2024LiSACReview} examine the integration of lighting, sensing, and
communication, with performance tables for visible-light and hybrid systems.
Lyu et al.\ \cite{Lyu2026PhotonicTHzWaveforms} review photonic millimeter-wave
and terahertz demonstrations and compare integrated waveforms through models
and experiments. Other focused reviews address optical transmission and
estimation, or communication and sensing in fiber networks
\cite{Su2025OISACTheoryPractice,Khorasgani2026GeneralOpticalISACSurvey,
Liu2026OpticalFiberISACReview,Lu2026ConvergedFiberISACReview}.

Broader surveys connect these technical developments to larger research
questions. Zhang et al.\ \cite{Zhang2026ISACEvolution} trace ISAC across
spectrum, network scale, and sensing modalities, including optical and
RF--optical systems. Mohsan et al.\ \cite{Mohsan2026OISACReview} bring together
optical platforms, experimental results, waveform design, and applications.
Both discuss performance measures and the need for further evaluation.
Together, these works explain design choices, report numerical
demonstrations, and identify research directions.

Table~\ref{tab:prior_surveys} distinguishes the physical coverage and review
approach of representative works. Its platform columns identify where the
sensing interaction takes place. The last two columns address how studies
are selected and how results from different studies are judged comparable.
This distinction matters because a controlled waveform comparison and a
comparison between independently reported experiments answer different
questions. The latter also requires agreement on what was measured and under
which conditions.

\begin{table*}[!t]
\caption{Physical coverage and review approach of representative related
reviews and tutorial-style syntheses}
\label{tab:prior_surveys}
\centering
\begingroup
\footnotesize
% Local Wasy glyphs keep the full and half-filled circles the same size.
\newcommand{\surveyGlyph}[1]{{\fontencoding{U}\fontfamily{wasy}\selectfont\char#1}}
\newcommand{\surveyFull}{\mbox{\surveyGlyph{32}}}
\newcommand{\surveyPartial}{\leavevmode
  \hbox to 0pt{\surveyGlyph{71}\hss}\hbox{\surveyGlyph{35}}}
\setlength{\tabcolsep}{3.5pt}
\renewcommand{\arraystretch}{1.2}
\renewcommand{\tabularxcolumn}[1]{m{#1}}
\begin{tabularx}{\textwidth}{@{}%
>{\raggedright\arraybackslash}m{0.13\textwidth}
>{\raggedright\arraybackslash}m{0.30\textwidth}
>{\centering\arraybackslash}m{0.055\textwidth}
>{\centering\arraybackslash}m{0.09\textwidth}
>{\centering\arraybackslash}m{0.09\textwidth}
>{\centering\arraybackslash}X
>{\centering\arraybackslash}X@{}}
\toprule
\textbf{Work / year} & \textbf{Main focus} &
\multicolumn{3}{c}{\textbf{Optical platforms}} &
\multicolumn{2}{c}{\textbf{Review approach}} \\
\cmidrule(lr){3-5}\cmidrule(l){6-7}
& & Fiber & Optical wireless & Photonic wireless &
Study selection & Cross-study comparison rules \\
\midrule
Wen et al.\ \cite{Wen2024OISACArchitectures}\newline 2024 &
Optical architectures, waveform design, and resource allocation &
\textemdash & \surveyFull & \textemdash & \textemdash & \textemdash \\
\addlinespace[3pt]
Liang et al.\ \cite{Liang2024LiSACReview}\newline 2024 &
Lighting--sensing--communication integration and hybrid links &
\textemdash & \surveyFull & \textemdash & \textemdash & \textemdash \\
\addlinespace[3pt]
Lyu et al.\ \cite{Lyu2026PhotonicTHzWaveforms}\newline 2026 &
Integrated photonic waveforms and controlled performance comparisons &
\textemdash & \textemdash & \surveyFull & \textemdash & \textemdash \\
\addlinespace[3pt]
Zhang et al.\ \cite{Zhang2026ISACEvolution}\newline 2026 &
ISAC evolution across spectrum, networks, and sensing modalities &
\surveyPartial & \surveyFull & \surveyFull & \textemdash & \surveyPartial \\
\addlinespace[3pt]
Mohsan et al.\ \cite{Mohsan2026OISACReview}\newline 2026 &
Optical platforms, experimental demonstrations, and applications &
\surveyFull & \surveyFull & \surveyFull & \textemdash & \surveyPartial \\
\midrule
\textbf{This survey} &
Shared-resource performance relations and evidence-based evaluation &
\surveyFull & \surveyFull & \surveyFull & \surveyFull & \surveyFull \\
\bottomrule
\end{tabularx}
\par\vspace{4pt}
\noindent
\parbox{\textwidth}{\raggedright\footnotesize
\surveyFull: structured treatment; \surveyPartial: selected discussion or an
identified research need; \textemdash: not developed in the stated sense.
Fiber denotes sensing in the fiber, not fiber transport alone; photonic
wireless denotes optically generated or processed signals with wireless
RF propagation. Only optical-related coverage is marked for Zhang et al.
Study selection denotes an explicit search and eligibility procedure.
Cross-study comparison rules address metric definitions, measurement points,
conditions, baselines, and supporting evidence; they are distinct from a
controlled comparison within one model or experiment.\par}
\endgroup
\end{table*}

A further need is to explain when a reported result can guide another
design. Zhang et al.\ identify the difficulty
of using common performance measures across RF and optical systems, while
Mohsan et al.\ call for shared benchmarks
\cite{Zhang2026ISACEvolution,Mohsan2026OISACReview}. We address this question
by linking the shared design choice to communication and sensing outcomes,
then examining the measurement conditions that support the link. This lets us
distinguish a result demonstrated in one setting from a design lesson that
can guide another. The same analysis identifies which joint measurement or
experiment should come next.

\subsection{Scope and Contributions}

This survey examines 206 unique O-ISAC studies first made publicly available
from 1 January 2020 through 22 June 2026. It covers guided optical systems,
optical wireless links, and photonics-enabled wireless systems, including
hybrid configurations. Across these settings, we examine how communication
and sensing are connected through a physical resource, a signal, or a joint
design objective. The review follows three questions. What is shared between
the two functions? How does the shared design affect their performance?
What further evaluation is needed to use the resulting insight in another
system?

The survey makes three contributions that connect system selection,
performance interpretation, and experimental evaluation.
\begin{itemize}
\item \textbf{A physical basis for selecting an O-ISAC design.}
We distinguish platforms by their sensing paths and identify how hardware,
signals, and allocated resources connect the two functions. This explains
which integration mechanisms fit a given physical setting and which
constraints must be considered when adapting them
(Sections~\ref{sec:foundations_comparison}
and~\ref{sec:optical_modality_families}).

\item \textbf{A common basis for interpreting performance.}
We relate bandwidth, power, waveform, and receiver choices to communication
rate and reliability alongside sensing accuracy and resolution. We distinguish
relationships demonstrated within a study from comparisons that require
compatible metrics, measurement points, and operating conditions. This
supports design lessons while keeping numerical results tied to the settings
in which they were obtained
(Sections~\ref{sec:foundations_comparison} and~\ref{sec:tradeoffs}).

\item \textbf{Evaluation priorities derived from the evidence.}
We connect validation findings to joint tests under traffic, changing
illumination, atmospheric variation, and optical drift. The proposed tests
pair communication and sensing measurements with separate-function baselines
and observable signs of progress. They give researchers concrete next steps
for testing whether a design remains useful beyond its original setting
(Sections~\ref{sec:validation_reproducibility}
and~\ref{sec:discussion_roadmap}).
\end{itemize}

Section~\ref{sec:foundations_comparison} introduces the physical and performance
concepts used in the review, and Section~\ref{sec:review_methods} describes the
review methods. Sections~\ref{sec:optical_modality_families}
through~\ref{sec:technologies_applications_6g} examine platforms, joint design,
validation, and applications. Section~\ref{sec:discussion_roadmap} develops the
research priorities, and Section~\ref{sec:conclusion} concludes the article.

\FloatBarrier
